\setcounter{section}{25}

  \zwischenueberschrift{Soal latihan}

\inputaufgabegibtloesung
{}
{

Misalkan $X$ suatu
\definitionsverweis {manifold diferensiabel}{}{}
dan
\maabb {p} { E  =  X \times \R } { X
} {}
bundel vektor trivial di atas $X$ yang ber-
\definitionsverweis {rank}{}{}
$1$. Misalkan $\omega$ suatu
$1$-\definitionsverweis {bentuk diferensial}{}{}
pada $X$. Tunjukkan pernyataan-pernyataan berikut.
\aufzaehlungfuenf{Pemetaan
\maabbeledisp {} {TE} { p^*E
} {(P,u,v,w) } { (P, u, \omega(P,v) +w)
} {,}
mendefinisikan suatu
\definitionsverweis {koneksi}{}{.}
}{Untuk fungsi-fungsi yang terdiferensialkan secara kontinu
\maabb {f} { U } {\R
} {}
pada suatu himpunan terbuka

\mavergleichskette
{\vergleichskette
{ U
}
{ \subseteq }{ X
}
{  }{
}
{    }{
}
{   }{
}
}
{}{}{}
turunan vertikal dari koneksi tersebut diberikan oleh

\mavergleichskettedisp
{\vergleichskette
{ \nabla (f) (P)
}
{ =} { \omega (P) + (df)_P
}
{ } {
}
{ } {
}
{ } {
}
}
{}{}{}
\zusatzklammer {di kedua ruas kita mengidentifikasi fungsi bernilai riil $f$ dengan penampang
\maabbele {} { X } { X \times \R
} { x } { (x,f(x))
} {.}} {} {}
}{Suatu fungsi yang terdiferensialkan secara kontinu
\maabb {f} { U } {\R
} {}
merupakan
\definitionsverweis {penampang horizontal}{}{}
di atas $U$ terhadap $\nabla$ jika dan hanya jika $-f$ merupakan
\definitionsverweis {bentuk primitif}{}{}
dari $\omega$ pada $U$.
}{Bentuk $\omega$
\definitionsverweis {tertutup}{}{}
jika dan hanya jika koneksi tersebut
\definitionsverweis {terintegralkan secara lokal}{}{.}
}{Bentuk $\omega$
\definitionsverweis {eksak}{}{}
jika dan hanya jika koneksi tersebut
\definitionsverweis {terintegralkan secara global}{}{.}
}

}
{} {}

\inputaufgabe
{}
{

Kita tinjau bundel vektor trivial
\maabbdisp {p} {E= \R^2 \times \R} { \R^2
} {.}
Deskripsikan suatu
\definitionsverweis {koneksi}{}{}
pada $E$ sedemikian sehingga tidak terdapat
\definitionsverweis {penampang horizontal}{}{.}

}
{} {}

\inputaufgabe
{}
{

Misalkan $X$ suatu
\definitionsverweis {manifold diferensiabel}{}{}
dan
\maabb {p} { E  =  X \times \R } { X
} {}
bundel vektor trivial di atas $X$ yang ber-
\definitionsverweis {rank}{}{}
$1$. Misalkan $\omega$ suatu
$1$-\definitionsverweis {bentuk diferensial}{}{}
pada $X$. Tunjukkan pernyataan-pernyataan berikut.
\aufzaehlungdrei{Pemetaan
\maabbeledisp {} {TE} { p^*E
} {(P,u,v,w) } { (P, u, u\omega(P,v) +w)
} {,}
mendefinisikan suatu
\definitionsverweis {koneksi}{}{.}
}{Untuk fungsi-fungsi yang terdiferensialkan secara kontinu
\maabb {f} { U } {\R
} {}
pada suatu himpunan terbuka

\mavergleichskette
{\vergleichskette
{ U
}
{ \subseteq }{ X
}
{  }{
}
{    }{
}
{   }{
}
}
{}{}{}
turunan vertikal dari koneksi tersebut diberikan oleh

\mavergleichskettedisp
{\vergleichskette
{ \nabla (f) (P)
}
{ =} { f\omega (P) + (df)_P
}
{ } {
}
{ } {
}
{ } {
}
}
{}{}{}
\zusatzklammer {di kedua ruas kita mengidentifikasi fungsi bernilai riil $f$ dengan penampang
\maabbele {} { X } { X \times \R
} { x } { (x,f(x))
} {.}} {} {}
}{Suatu fungsi yang terdiferensialkan secara kontinu dan tidak pernah nol
\maabb {f} { U } {\R
} {}
merupakan
\definitionsverweis {penampang horizontal}{}{}
di atas $U$ terhadap $\nabla$ jika dan hanya jika $-   \ln  \betrag {  f  }$ merupakan
\definitionsverweis {bentuk primitif}{}{}
dari $\omega$ pada $U$.
}

}
{} {}

\inputaufgabe
{}
{

Misalkan $X$ suatu
\definitionsverweis {manifold diferensiabel}{}{}
dan $E$ suatu
\definitionsverweis {bundel vektor diferensiabel}{}{}
di atas $X$ yang dilengkapi dengan suatu
\definitionsverweis {koneksi}{}{.}
Tunjukkan bahwa suatu penampang yang terdiferensialkan secara kontinu
\maabb {s} { U } { E\vert_U
} {}
pada suatu himpunan terbuka

\mavergleichskette
{\vergleichskette
{ U
}
{ \subseteq }{ X
}
{  }{
}
{    }{
}
{   }{
}
}
{}{}{}
bersifat
\definitionsverweis {horizontal}{}{}
jika dan hanya jika
\definitionsverweis {turunan vertikal}{}{}-nya memenuhi

\mavergleichskette
{\vergleichskette
{ \nabla s
}
{ = }{ 0
}
{  }{
}
{    }{
}
{   }{
}
}
{}{}{.}

}
{} {}

\inputaufgabe
{}
{

Misalkan $X$ suatu
\definitionsverweis {manifold diferensiabel}{}{}
dan $E$ suatu
\definitionsverweis {bundel vektor diferensiabel}{}{}
di atas $X$ yang dilengkapi dengan suatu
\definitionsverweis {koneksi}{}{.}
Misalkan $Z$ suatu manifold diferensiabel lain dengan pemetaan diferensiabel
\maabb {\psi} { Z } { X
} {.}
Tunjukkan bahwa suatu penampang yang terdiferensialkan secara kontinu
\maabb {s} { Z } { E
} {}
di atas $\psi$ bersifat
\definitionsverweis {horizontal}{}{}
jika dan hanya jika
\definitionsverweis {turunan vertikal}{}{}-nya memenuhi

\mavergleichskette
{\vergleichskette
{ \nabla s
}
{ = }{ 0
}
{  }{
}
{    }{
}
{   }{
}
}
{}{}{.}

}
{} {}

Untuk soal berikut, gunakan
Soal 25.21.

\inputaufgabe
{}
{

Misalkan
\maabb {p} { E } { X
} {}
suatu
\definitionsverweis {bundel vektor}{}{}
ber-
\definitionsverweis {rank}{}{}
$r$ di atas suatu
\definitionsverweis {manifold diferensiabel}{}{}
$X$, yang dilengkapi dengan suatu
\definitionsverweis {koneksi yang terintegralkan secara global}{}{.}
Tunjukkan bahwa terdapat
\zusatzklammer {yang terdefinisi pada $X$} {} {}
\definitionsverweis {penampang-penampang horizontal}{}{}

\mathl{s_1  , \ldots , s_r}{} di $E$ sedemikian sehingga
\maabbeledisp {} {X \times \R^r} { E
} {(P,v_1  , \ldots , v_r)} { \sum_{i = 1}^r v_i s_i(P)
} {,}
merupakan suatu
\definitionsverweis {isomorfisme bundel vektor}{}{.}

}
{} {}

\inputaufgabegibtloesung
{}
{

Misalkan
\maabb {p} { E } { X
} {}
suatu
\definitionsverweis {bundel vektor}{}{}
ber-
\definitionsverweis {rank}{}{}
$r$ di atas suatu
\definitionsverweis {manifold diferensiabel}{}{}
$X$, yang dilengkapi dengan suatu koneksi linear yang
\definitionsverweis {terintegralkan secara global}{}{.}
Tunjukkan bahwa terdapat
\zusatzklammer {yang terdefinisi pada $X$} {} {}
\definitionsverweis {penampang-penampang horizontal}{}{}

\mathl{s_1 , \ldots , s_r}{} di $E$ sedemikian sehingga, di bawah
\definitionsverweis {isomorfisme bundel vektor}{}{}
\maabbeledisp {} {X \times \R^r} { E
} {(P,v_1 , \ldots , v_r)} { \sum_{i = 1}^r v_i s_i(P)
} {,}
\definitionsverweis {simbol-simbol Christoffel}{}{}
dari koneksi tersebut terhadap penampang basis
\mathl{s_1  , \ldots , s_r}{} semuanya bernilai nol.

}
{} {}

\inputaufgabegibtloesung
{}
{

Misalkan diberikan suatu
\definitionsverweis {medan vektor}{}{}
kontinu yang bergantung pada waktu,
\maabbeledisp {F} {I \times \R^n } { \R^n
} { \left(  t   , \,  u_1  , \,   \ldots , \,  u_n                \right) } { \left(  F_1 \left(  u_1 , \ldots ,  u_n  \right)    , \,   \ldots  , \,   F_n \left(  u_1 , \ldots , u_n  \right)                 \right)
} {,}
beserta
\definitionsverweis {sistem persamaan diferensial biasa}{}{}
yang terkait,

\mavergleichskettedisp
{\vergleichskette
{ u'
}
{ =} { F(t,u)
}
{ } {
}
{ } {
}
{ } {
}
}
{}{}{}
pada suatu
\definitionsverweis {interval terbuka}{}{}

\mavergleichskette
{\vergleichskette
{ I
}
{ \subseteq }{ \R
}
{  }{
}
{    }{
}
{   }{
}
}
{}{}{.}
\aufzaehlungdrei{Tunjukkan bahwa pemetaan
\maabbeledisp {} {I \times \R^n \times \R \times \R^n } { I \times \R^n  \times \R^n
} {(t,u,v,w)} { \left(  t   , \,  u  , \,   - vF \left(  t, u_1  , \ldots , u_n  \right)  +w                 \right)
} {,}
memberikan
\zusatzklammer {proyeksi vertikal dari suatu} {} {}
\definitionsverweis {koneksi}{}{}
pada bundel vektor trivial
\maabbdisp {} {I \times \R^n } { \R^n
} {}
diberikan.
}{Tentukan
\definitionsverweis {turunan vertikal}{}{}
$\nabla_\partial u$ dari suatu
\definitionsverweis {kurva diferensiabel}{}{}
\maabbdisp {u} { I } { \R^n
} {}
\zusatzklammer {yang dipandang sebagai penampang di $I \times \R^n$} {} {.}
}{Tunjukkan bahwa suatu kurva diferensiabel
\maabb {u} { I } { \R^n
} {}
merupakan solusi dari sistem persamaan diferensial jika dan hanya jika $u$ merupakan
\definitionsverweis {penampang horizontal}{}{}
terhadap koneksi tersebut.
}

}
{} {}

\inputaufgabe
{}
{

Misalkan $\nabla$ suatu
\definitionsverweis {koneksi linear}{}{}
pada suatu
\definitionsverweis {bundel vektor diferensiabel}{}{}
$E$ di atas suatu
\definitionsverweis {manifold diferensiabel}{}{}
$X$. Tunjukkan bahwa penampang nol bersifat
\definitionsverweis {horizontal}{}{.}

}
{} {}

\inputaufgabe
{}
{

Misalkan $\nabla$
\definitionsverweis {koneksi trivial}{}{}
pada bundel vektor trivial
\maabb {p} {X \times W } { X
} {}
yang terkait dengan suatu ruang vektor riil
\mathl{W}{} di atas suatu
\definitionsverweis {manifold diferensiabel}{}{}
$X$. Tunjukkan bahwa $\nabla$
\definitionsverweis {linear}{}{.}

}
{} {}

\inputaufgabegibtloesung
{}
{

Misalkan

\mavergleichskette
{\vergleichskette
{ U
}
{ \subseteq }{ \R^n
}
{  }{
}
{    }{
}
{   }{
}
}
{}{}{}
suatu
\definitionsverweis {himpunan bagian terbuka}{}{}
dan misalkan $\nabla$
\definitionsverweis {koneksi trivial}{}{}
pada
\mathl{U \times \R}{.} Tunjukkan bahwa
\definitionsverweis {simbol-simbol Christoffel}{}{}
dan
\definitionsverweis {turunan vertikal}{}{}
dari koneksi tersebut memenuhi sifat-sifat berikut.
\aufzaehlungvier{Simbol-simbol Christoffel memenuhi

\mavergleichskettedisp
{\vergleichskette
{ \Gamma_{ij}^k
}
{ =} { 0
}
{ } {
}
{ } {
}
{ } {
}
}
{}{}{}
untuk semua $i,j,k$.
}{Berlaku

\mavergleichskettedisp
{\vergleichskette
{ \nabla_{\partial_i }  \partial_j
}
{ =} { 0
}
{ } {
}
{ } {
}
{ } {
}
}
{}{}{}
untuk
\definitionsverweis {medan-medan vektor standar}{}{}
$\partial_i$.
}{Berlaku

\mavergleichskettedisp
{\vergleichskette
{ \nabla_{\partial_i }  \left(  \sum_{j = 1}^n f_j \partial_j  \right)
}
{ =} { \sum_{j = 1}^n \partial_i(f_j) \partial_j
}
{ } {
}
{ } {
}
{ } {
}
}
{}{}{}
untuk fungsi-fungsi diferensiabel $f_j$.
}{Untuk suatu medan vektor $V$ dan suatu medan vektor diferensiabel $W$ pada $U$ berlaku

\mavergleichskettedisp
{\vergleichskette
{ ( \nabla_V W)(P)
}
{ =} { \left(  D W   \right)_{ P } \left(   V(P)  \right)
}
{ } {
}
{ } {
}
{ } {
}
}
{}{}{.}
}

}
{} {}

\inputaufgabegibtloesung
{}
{

Pada
\definitionsverweis {bundel vektor trivial}{}{}

\mathl{\R \times \R^2}{} di atas $\R$, kita tinjau
\definitionsverweis {koneksi trivial}{}{}
$\nabla$.
\aufzaehlungdrei{Tunjukkan bahwa

\mavergleichskette
{\vergleichskette
{ f_1(t)
}
{ = }{ \left(  1+ t^2  , \,  t^3                   \right)
}
{  }{
}
{    }{
}
{   }{
}
}
{}{}{}
dan

\mavergleichskette
{\vergleichskette
{ f_2(t)
}
{ = }{ \left(  t  , \,   1+t^2                   \right)
}
{  }{
}
{    }{
}
{   }{
}
}
{}{}{}
merupakan
\definitionsverweis {penampang-penampang basis}{}{}
dari bundel vektor tersebut.
}{Nyatakan penampang basis standar $e_1,e_2$ sebagai kombinasi linear dari penampang basis $f_1,f_2$.
}{Tentukan
\definitionsverweis {simbol-simbol Christoffel}{}{}
dari $\nabla$ terhadap penampang basis tersebut.
}

}
{} {}

\inputaufgabe
{}
{

Pada
\definitionsverweis {bundel vektor trivial}{}{}

\mathl{\R^2 \times \R^2}{} di atas $\R^2$, kita tinjau
\definitionsverweis {koneksi trivial}{}{}
$\nabla$.
\aufzaehlungzwei {Tunjukkan bahwa

\mavergleichskette
{\vergleichskette
{ f_1(x,y)
}
{ = }{ \left(  x   , \,   y                   \right)
}
{  }{
}
{    }{
}
{   }{
}
}
{}{}{}
dan

\mavergleichskette
{\vergleichskette
{ f_2(x,y)
}
{ = }{ \left(  - y  , \,   x                   \right)
}
{  }{
}
{    }{
}
{   }{
}
}
{}{}{}
merupakan
\definitionsverweis {penampang-penampang basis}{}{}
dari bundel vektor pada

\mavergleichskette
{\vergleichskette
{ U
}
{ = }{ \R^2 \setminus \{ (0,0)\}
}
{  }{
}
{    }{
}
{   }{
}
}
{}{}{.}
} {Tentukan
\definitionsverweis {simbol-simbol Christoffel}{}{}
dari $\nabla$ terhadap penampang basis tersebut pada $U$.
}

}
{} {}

\inputaufgabegibtloesung
{}
{

Misalkan diberikan suatu
\definitionsverweis {sistem persamaan diferensial linear homogen}{}{}

\mavergleichskettedisp
{\vergleichskette
{ u'
}
{ =} { Mu
}
{ } {
}
{ } {
}
{ } {
}
}
{}{}{}
pada suatu
\definitionsverweis {interval terbuka}{}{}

\mavergleichskette
{\vergleichskette
{ I
}
{ \subseteq }{ \R
}
{  }{
}
{    }{
}
{   }{
}
}
{}{}{}
dengan

\mavergleichskettedisp
{\vergleichskette
{ M(t)
}
{ =} { \left(  a_{kj} (t)  \right)_{kj}
}
{ } {
}
{ } {
}
{ } {
}
}
{}{}{}
\zusatzklammer {di sini $k$ adalah indeks baris dan $j$ adalah indeks kolom} {} {}
serta fungsi-fungsi kontinu
\maabbdisp {a_{kj}} { I } {\R
} {}
diberikan.
Melalui persamaan

\mavergleichskettedisp
{\vergleichskette
{ \Gamma_{j}^k (t)
}
{ \defeq} { -a_{kj}
}
{ } {
}
{ } {
}
{ } {
}
}
{}{}{}
kita memandang fungsi-fungsi ini sebagai
\definitionsverweis {simbol-simbol Christoffel}{}{}
\zusatzklammer {dengan

\mavergleichskettek
{\vergleichskettek
{ i
}
{ = }{ 1
}
{  }{
}
{    }{
}
{   }{
}
}
{}{}{}} {} {}
dari
\definitionsverweis {koneksi linear}{}{}
$\nabla$ pada bundel vektor trivial
\maabb {p} {I \times \R^n } { I
} {}
dalam pengertian
Catatan 25.8.
Tunjukkan bahwa suatu
\definitionsverweis {kurva diferensiabel}{}{}
\maabbdisp {u} { I } { \R^n
} {}
merupakan solusi sistem persamaan diferensial jika dan hanya jika $u$ merupakan
\definitionsverweis {penampang horizontal}{}{}
terhadap koneksi tersebut.

}
{} {}

\inputaufgabe
{}
{

Misalkan

\mavergleichskette
{\vergleichskette
{ Y
}
{ \subseteq }{ W
}
{ \subseteq }{ \R^n
}
{    }{
}
{   }{
}
}
{}{}{,}
$W$
\definitionsverweis {terbuka}{}{,}
suatu
\definitionsverweis {hipermuka terdiferensialkan}{}{,}
dan
\maabb {\gamma} { I } { Y
} {}
suatu
\definitionsverweis {kurva diferensiabel}{}{.}
Misalkan
\maabbdisp {F} { I } { TY
} {}
suatu medan vektor diferensiabel sepanjang $\gamma$; dengan kata lain, terdapat diagram komutatif

\mathdisp {\begin{matrix}  &   &   & TY  \\ & & F \nearrow & \downarrow  \\  &  I  & \stackrel{ \gamma }{\longrightarrow}  &  Y \! \\ \end{matrix}} {  }

Tunjukkan bahwa $F$
\definitionsverweis {paralel sepanjang}{}{}
$\gamma$ jika dan hanya jika $F$ merupakan
\definitionsverweis {penampang horizontal}{}{}
terhadap
\zusatzklammer {yang didefinisikan menggunakan proyeksi ortogonal} {} {}
\definitionsverweis {koneksi}{}{}
pada $TY$.

}
{} {}

\inputaufgabe
{}
{

Misalkan

\mavergleichskette
{\vergleichskette
{ Y
}
{ \subseteq }{ W
}
{ \subseteq }{ \R^n
}
{    }{
}
{   }{
}
}
{}{}{,}
$W$
\definitionsverweis {terbuka}{}{,}
suatu
\definitionsverweis {hipermuka terdiferensialkan}{}{,}
dan misalkan $\nabla$
\zusatzklammer {yang didefinisikan menggunakan proyeksi ortogonal} {} {}
\definitionsverweis {koneksi}{}{}
pada $TY$. Tunjukkan bahwa $\nabla$
\definitionsverweis {linear}{}{.}

}
{} {}

\inputaufgabe
{}
{

Misalkan
\mathkor {} {\nabla_1} {dan} {\nabla_2} {}
\definitionsverweis {koneksi-koneksi linear}{}{}
pada suatu
\definitionsverweis {bundel vektor diferensiabel}{}{}
$E$ di atas suatu
\definitionsverweis {manifold diferensiabel}{}{}
$X$. Misalkan
\maabbdisp {g_1,g_2} { X } { \R
} {}
fungsi-fungsi yang
\definitionsverweis {terdiferensialkan secara kontinu}{}{}
dengan

\mavergleichskette
{\vergleichskette
{ g_1+g_2
}
{ = }{ 1
}
{  }{
}
{    }{
}
{   }{
}
}
{}{}{.}
Tunjukkan bahwa
\mathl{g_1 \nabla_1+ g_2 \nabla_2}{} juga merupakan suatu koneksi linear
\zusatzklammer {dalam pengertian yang telah didefinisikan} {?} {}
pada $E$.

}
{} {}

\inputaufgabe
{}
{

Misalkan

\mavergleichskette
{\vergleichskette
{ X
}
{ = }{ \bigcup_{i \in I} U_i
}
{  }{
}
{    }{
}
{   }{
}
}
{}{}{}
suatu
\definitionsverweis {tutupan terbuka}{}{}
dari suatu
\definitionsverweis {manifold diferensiabel}{}{}
$X$ yang mentrivialkan
\definitionsverweis {bundel vektor diferensiabel}{}{}
$E$. Misalkan $\nabla_i$
\definitionsverweis {koneksi trivial}{}{}
pada $E\vert_{U_i }$ dan misalkan
\mathkor {} {h_j} {} {j \in J} {,}
suatu
\definitionsverweis {partisi satuan}{}{}
yang subordinat terhadap tutupan tersebut,

\mathbed {h_j} {}
 {j \in J} {}
 {} {} {} {.}
Tunjukkan bahwa $\sum_{j \in J} h_j \nabla_{i(j)}$ merupakan suatu koneksi linear yang terdefinisi dengan baik pada $E$.

}
{} {}

\inputaufgabe
{}
{

Misalkan $X$ suatu
\definitionsverweis {manifold diferensiabel}{}{}
dan $\omega$ suatu
$1$-\definitionsverweis {bentuk diferensial}{}{}
kontinu pada $X$. Tunjukkan bahwa pemetaan
\maabbeledisp {} {C^1(X,\R)} { C^0(X,T^*X)
} { f } { \nabla_\omega (f) \defeq f \omega + df
} {,}
memenuhi aturan Leibniz dari
Teorema 25.5.

}
{} {}

\inputaufgabe
{}
{

Misalkan
\maabb {p} { E } { X
} {}
suatu
\definitionsverweis {bundel vektor}{}{}
diferensiabel di atas suatu
\definitionsverweis {manifold diferensiabel}{}{}
$X$. Misalkan $E$ sendiri merupakan suatu
\definitionsverweis {manifold Riemann}{}{.}
Tunjukkan bahwa suatu
\definitionsverweis {koneksi linear}{}{}
pada $E$ diperoleh dengan mengambil, untuk
\definitionsverweis {bundel vertikal}{}{}

\mavergleichskette
{\vergleichskette
{ V
}
{ \subseteq }{ T E
}
{  }{
}
{    }{
}
{   }{
}
}
{}{}{}
\definitionsverweis {komplemen ortogonal}{}{}
sebagai suatu
\definitionsverweis {bundel horizontal}{}{.}

}
{} {}

\inputaufgabe
{}
{

Misalkan
\maabb {p} { E } { X
} {}
suatu
\definitionsverweis {bundel vektor}{}{}
di atas
\definitionsverweis {manifold diferensiabel}{}{}
$X$ dengan
\definitionsverweis {basis topologi terhitung}{}{.}
Gunakan
Soal 22.15
dan
Soal 25.18
untuk menunjukkan bahwa terdapat suatu
\definitionsverweis {koneksi linear}{}{}
pada $E$.

}
{} {}

  \zwischenueberschrift{Soal yang dikumpulkan}

\inputaufgabe
{5 (3+2)}
{

Misalkan

\mavergleichskette
{\vergleichskette
{ X
}
{ = }{ S^1
}
{  }{
}
{    }{
}
{   }{
}
}
{}{}{}
\definitionsverweis {lingkaran satuan}{}{}
dengan
\definitionsverweis {bundel tangen}{}{}
\maabb {p} {TS^1 = S^1 \times \R} { S^1
} {.}
\aufzaehlungzwei {Deskripsikan suatu
\definitionsverweis {koneksi}{}{}
pada $TS^1$ sedemikian sehingga tidak terdapat
\definitionsverweis {penampang horizontal}{}{}
pada seluruh $S^1$.
} {Tentukan suatu penampang horizontal dari koneksi pada (1) sepanjang pemetaan
\maabbeledisp {} {\R} { S^1
} { t } { \left(  \cos  t   , \,   \sin  t                   \right)
} {.}
}

}
{} {}

\inputaufgabe
{4}
{

Misalkan $X$ suatu
\definitionsverweis {manifold diferensiabel}{}{}
dan $E$ suatu
\definitionsverweis {bundel vektor diferensiabel}{}{}
di atas $X$ yang dilengkapi dengan suatu
\definitionsverweis {koneksi}{}{.}
Misalkan $Z$ suatu manifold diferensiabel lain dengan pemetaan diferensiabel
\maabb {\psi} { Z } { X
} {.}
Misalkan $Z$
\definitionsverweis {terhubung}{}{}
dan misalkan
\maabb {s_1,s_2} { Z } { E
} {}
\definitionsverweis {penampang-penampang horizontal}{}{}
di atas $\psi$ dengan

\mavergleichskette
{\vergleichskette
{ s_1(z)
}
{ = }{ s_2(z)
}
{  }{
}
{    }{
}
{   }{
}
}
{}{}{}
untuk suatu titik

\mavergleichskette
{\vergleichskette
{ z
}
{ \in }{ Z
}
{  }{
}
{    }{
}
{   }{
}
}
{}{}{.}
Tunjukkan bahwa

\mavergleichskette
{\vergleichskette
{ s_1
}
{ = }{ s_2
}
{  }{
}
{    }{
}
{   }{
}
}
{}{}{.}

}
{} {}

\inputaufgabe
{3}
{

Misalkan $X$ suatu
\definitionsverweis {manifold diferensiabel}{}{}
dan $E$ suatu
\definitionsverweis {bundel vektor diferensiabel}{}{}
di atas $X$ yang dilengkapi dengan suatu
\definitionsverweis {koneksi linear}{}{,}
dan misalkan $G$ suatu
\definitionsverweis {medan vektor}{}{}
pada $X$. Tunjukkan bahwa untuk setiap
\definitionsverweis {penampang diferensiabel}{}{}
$s$ di $E$ dan setiap
\definitionsverweis {fungsi diferensiabel}{}{}
$f$
\zusatzklammer {yang keduanya terdefinisi pada suatu himpunan terbuka

\mavergleichskette
{\vergleichskette
{ U
}
{ \subseteq }{ X
}
{  }{
}
{    }{
}
{   }{
}
}
{}{}{}} {} {}
berlaku hubungan

\mavergleichskettedisp
{\vergleichskette
{ \nabla_G (fs)
}
{ =} { s \otimes D_G (f) + f \nabla_G (s)
}
{ } {
}
{ } {
}
{ } {
}
}
{}{}{.}

}
{} {}

\inputaufgabe
{3}
{

Berikan suatu contoh
\definitionsverweis {koneksi linear}{}{}
pada suatu
\definitionsverweis {bundel vektor}{}{}
\maabb {p} { E } { X
} {}
di atas suatu
\definitionsverweis {manifold diferensiabel}{}{}
$X$ sedemikian sehingga ruang vektor
\mathl{H(\nabla,X)}{} dari
\definitionsverweis {penampang-penampang horizontal}{}{}
memiliki
\definitionsverweis {dimensi ruang vektor}{}{}
yang lebih besar daripada
\definitionsverweis {rank}{}{}
$E$.

}
{} {}
